\begin{figure}[H]
\centering
\begin{tikzpicture}[>=Stealth,x=1cm,y=1cm,font=\small]
\fill[black!6] (-4,-.8) rectangle (-1,.8);
\fill[black!6] (1,-.8) rectangle (4,.8);
\draw[->] (-4.5,0)--(4.7,0) node[right] {$\RealPart\ResolventParameter$};
\draw[->] (0,-1.05)--(0,1.3) node[above] {$\ImaginaryPart\ResolventParameter$};
\draw[thick] (-4,-.8) rectangle (-1,.8);
\draw[thick] (1,-.8) rectangle (4,.8);
\draw[->,thick] (-3.2,-.8)--(-2.5,-.8);
\draw[->,thick] (2,-.8)--(2.7,-.8);
\foreach \x in {-3.5,-2.8,-1.7,1.6,2.5,3.4} \fill (\x,0) circle (1.5pt);
\foreach \x in {-.55,-.3,.25,.5} \fill[black!45] (\x,0) circle (1.2pt);
\node[below left] at (0,0) {$0$};
\node[below] at (-4,-.85) {$-\SpectralOuterBound$};
\node[below] at (-1,-.85) {$-\SpectralCutoff$};
\node[below] at (1,-.85) {$\SpectralCutoff$};
\node[below] at (4,-.85) {$\SpectralOuterBound$};
\node[above] at (-2.5,.85) {$\SpectralContour$};
\node[above] at (2.5,.85) {$\SpectralContour$};
\end{tikzpicture}
\caption{The positively oriented boundary
$\SpectralContour=\partial\SpectralRegion$
encloses the retained positive and negative spectral values.
The dots are schematic.
The gaps at
$\pm\SpectralCutoff$
allow the same boundary to be used for
$\AmbientOperator_\SequenceIndex$
for all sufficiently large
$\SequenceIndex$.}
\label{fig:spectral-contour}
\end{figure}
